\section{Introduction}
\label{sec:intro}

The rapid proliferation of deep neural networks has led to an abundance of specialized models, each fine-tuned to excel on a particular target task or domain. However, maintaining and deploying a distinct, full-parameter model for every individual task is prohibitively expensive in terms of storage, memory footprint, and deployment latency. To circumvent this challenge, the machine learning community has turned to parameter-efficient multi-task learning techniques, among which \emph{model merging} has emerged as an exceptionally elegant paradigm \cite{wortsman2022robust, ilharco2022editing, tiesmerging}. By linearly interpolating the weights of multiple task-specific experts fine-tuned from a shared pretrained base model, model merging enables the creation of a single, unified multi-task model without any additional training parameters.

While static model merging methods---such as Task Arithmetic \cite{ilharco2022editing} or TIES-Merging \cite{tiesmerging}---rely on fixed coefficients to combine task experts, they are fundamentally limited when dealing with highly heterogeneous task mixtures. Under such conditions, a single static set of coefficients often fails to balance the trade-offs between disparate domains, leading to interference and sub-optimal performance. To address this, \emph{dynamic model merging} has been proposed, where a lightweight routing head dynamically determines sample-specific or batch-specific merging coefficients at test time \cite{adamerging}. This dynamic routing allows the merged model to adapt its internal representations on the fly, offering a highly flexible and parameter-efficient multi-task solution.

Despite these advantages, calibrating dynamic routing heads presents severe generalization challenges, particularly in low-data regimes. In practical scenarios, we often only have access to a tiny validation or calibration set (e.g., 16 samples per task) to optimize the routing head parameters. Under such extreme constraints, standard unregularized dynamic routers are highly susceptible to catastrophic overfitting, or worse, \emph{representational collapse} on high-conflict datasets. For instance, when merging experts specialized in highly distinct domains like handwritten digits (MNIST) and complex natural images (CIFAR-10), the routing head's optimization landscape becomes extremely volatile, frequently collapsing the performance of difficult or out-of-domain tasks like SVHN.

We argue that progress in model merging is best achieved through rigorous, scale-controlled baseline optimization, exhaustive sweeps, and robust ablation studies rather than unnecessarily complex theoretical metaphors. In this work, we deconstruct previous routing failures and demonstrate that they are largely an artifact of unregularized baseline optimization and the inherent limitations of the Softmax zero-sum competitive bottleneck. By replacing the Softmax router with a decoupled, independent \emph{Bounded Sigmoidal Router (BSigmoid-Router)} and introducing a simple yet highly principled regularizer, we can dramatically stabilize and enhance multi-task calibration.

We propose \textbf{Task-Correlation Prior Regularization (TCPR)}, a novel regularization technique that injects a pre-computed cross-task similarity matrix $S \in \mathbb{R}^{K \times K}$ as a prior to guide the optimization of routing projection weights during low-data calibration. We formulate two distinct variants of this similarity prior:
\begin{enumerate}
    \item \textbf{Parameter-Space Similarity (TCPR-Param):} The average cosine similarity of the task vectors across all layers of the network.
    \item \textbf{Representation-Space Similarity (TCPR-Rep):} The cosine similarity of intermediate representations extracted from the base model on a generic validation subset.
\end{enumerate}
Using the similarity matrix $S$, TCPR penalizes diverging routing signatures for highly correlated tasks while enforcing orthogonal routing paths for conflicting tasks. This directly resolves the competitive bottleneck during calibration, enabling stable convergence and preventing performance collapse on difficult tasks.

To validate our proposed approach, we conduct an exceptionally rigorous and exhaustive empirical evaluation. Using a compact Vision Transformer (\texttt{vit\_tiny\_patch16\_224}) backbone, we trainconverged task experts across four highly diverse datasets (MNIST, FashionMNIST, CIFAR-10, and SVHN) and calibrate the routing heads on an extremely challenging 64-sample set. Through a comprehensive logarithmic hyperparameter sweep over $\beta \in [10^{-6}, 10^2]$, we demonstrate that TCPR consistently outperforms standard L2-regularized baselines and state-of-the-art wave-interference methods like QWS-Merge. 

Our main contributions are summarized as follows:
\begin{itemize}
    \item We identify and analyze the representational collapse of classical routing heads on high-conflict datasets under low-data calibration, demonstrating that standard unregularized optimization is the primary root cause.
    \item We propose \textbf{Task-Correlation Prior Regularization (TCPR)}, introducing two distinct variants (TCPR-Param and TCPR-Rep) to incorporate task-relatedness priors into dynamic routing head calibration.
    \item We evaluate our methods against a highly rigorous set of seven baselines across a heterogeneous four-task Vision Transformer benchmark, showing that TCPR achieves state-of-the-art joint multi-task performance.
    \item We present exhaustive hyperparameter sensitivity sweeps and ablation analyses, establishing a robust, scale-invariant, and highly practical guide for low-data model merging.
\end{itemize}
